\documentclass[11pt]{article}

\usepackage[a4paper,margin=25mm]{geometry}
\usepackage{amsmath,amssymb}
\usepackage{booktabs}
\usepackage{graphicx}
\usepackage{tikz}
\usepackage[colorlinks=true,linkcolor=blue,citecolor=blue]{hyperref}

\title{TC23B: Crack Growth in a Skin-Bridged Stringer at a Fastener Hole\\
\large Implementation Notes and Benchmark Reconciliation for SCT-CG1}
\author{Southward Consulting Ltd}
\date{June 2026}

\newcommand{\ksi}{\,\mathrm{ksi}}
\newcommand{\inch}{\,\mathrm{in}}
\newcommand{\ksqin}{\,\mathrm{ksi\sqrt{in}}}

\begin{document}
\maketitle

\begin{abstract}
This note documents the TC23B crack case implemented in the SCT-CG1 crack
growth tool: NASGRO crack case TC23 (unequal through cracks at an offset hole
in a finite-width strip) extended with (i) selectable bearing-load $\beta$
models, and (ii) crack bridging by a fuselage skin attached through a single
vertical line of fasteners passing through the critical hole, solved by the
Swift displacement-compatibility method with fastener stiffness from the
modified Tate \& Rosenfeld equation. Verification against first-principles
limit cases and reconciliation against an independent benchmark code are
summarised, and the open questions for discussion with the benchmark author
are identified. Basic LEFM and NASGRO theory are assumed known and are not
repeated.
\end{abstract}

\section{Scope and configuration}

The intended application is fatigue cracking in a stringer, modelled by its
unfolded width $W$, containing a fastener hole of diameter $D$ at distance
$m$ from one edge, with through cracks $c_1$ (wide side) and $c_2$ (edge
side) growing from the hole. The stringer is attached to a background
\emph{infinite} skin by a single vertical row of fasteners at pitch $p$
passing through the critical hole (Fig.~\ref{fig:config}). As the crack
opens, the fasteners above and below the crack plane transfer load into the
intact skin, restraining the opening displacement and retarding growth.

\begin{figure}[h]
\centering
\begin{tikzpicture}[scale=0.9]
  % stringer strip
  \draw[thick] (0,0) rectangle (3.0,6.0);
  % fastener line
  \draw[dashed,gray] (2.0,0.2) -- (2.0,5.8);
  % fasteners
  \foreach \y in {1.0,1.8,4.2,5.0} \fill[blue!60] (2.0,\y) circle (0.07);
  % hole + cracks
  \draw[thick] (2.0,3.0) circle (0.22);
  \draw[red,very thick] (0.6,3.0) -- (1.78,3.0);
  \draw[red,very thick] (2.22,3.0) -- (2.75,3.0);
  % dimensions
  \draw[<->] (2.0,-0.4) -- (3.0,-0.4) node[midway,below] {$m$};
  \draw[<->] (0,-0.9) -- (3.0,-0.9) node[midway,below] {$W$};
  \draw[<->] (2.25,4.2) -- (2.25,5.0) node[midway,right] {$p$};
  \node[red] at (1.2,3.25) {$c_1$};
  \node[red] at (2.5,3.3) {$c_2$};
  % load arrows
  \foreach \x in {0.5,1.5,2.5} {
    \draw[->] (\x,6.1) -- (\x,6.6);
    \draw[->] (\x,-1.2) -- (\x,-1.7);
  }
  \node at (3.5,6.4) {$S_0$};
  % skin note
  \node[blue!60,anchor=west] at (3.4,3.0)
    {\footnotesize fastener row attached to infinite skin ($E_2,t_2,\nu_2$)};
\end{tikzpicture}
\caption{TC23B configuration. The crack plane is normal to the load; the
fastener row is collinear with the load through the critical hole.}
\label{fig:config}
\end{figure}

\section{Baseline TC23 implementation}

The unbridged solution follows NASGRO~8.2 Appendix~C, crack case TC23,
verified term-by-term against the manual:

\begin{itemize}
\item \textbf{Solution A} (remote tension): $\beta_i^A=\beta_i^{A1}\beta_i^{A2}\beta_i^{A3}$
  with the Tweed--Rooke/Bowie fit for unequal cracks at a hole in an
  infinite plate ($\beta_r$ polynomial and $\beta_u$ rational fit) and the
  $\varphi_1/\varphi_2$ finite-width corrections for crack-edge (A2) and
  hole-edge (A3) interaction.
\item \textbf{Solution B} (in-plane bending): $\beta^{B1}$ compounding with the
  $\Phi_1(d_1,d_2)$ FEM benchmark table (pp.~C-40/41).
\item \textbf{Solution C} (pin load): $\beta_i^C=\tfrac12(\beta_i^{C1}+\beta_i^{C2}+\beta_i^{C3})$,
  $K_i=\bigl(\beta_i^A S_0+\beta_i^B S_2+\tfrac{D}{W}\beta_i^C \bar S_3\bigr)\sqrt{\pi c_i}$,
  with eccentricity convention $e=W/2-B$ per the manual. Alternative bearing
  models are discussed in \S\ref{sec:bearing}.
\item \textbf{Growth law}: full NASGRO equation (Newman closure, $C_{th}^{\pm}$
  threshold fanning, thickness-corrected $K_c$); net-section yield on flow
  stress $(\sigma_{ys}+\sigma_u)/2$; optional every-cycle Irwin plastic-zone
  correction (two fixed-point iterations on $c+r_y$).
\item \textbf{Post-link-up}: when the edge-side ligament fails the analysis
  continues as a single edge crack (SENT $\beta$, $\eta$-interpolated, with
  the Gross--Srawley bending factor for $S_2$) rather than terminating.
\end{itemize}

Dual tips are integrated independently with a shared $\Delta N$ per step and
adaptive $\Delta c$ keyed to the fracture margin.

\section{Benchmark reconciliation}
\label{sec:bench}

Two cases were compared against an independent benchmark code
(V1Aerospace, v1.2.3): $D=0.2$, $t=0.05$, $c_{1,0}=0.01$, $c_{2,0}=0.05$,
$S_0=15\ksi$, $R=0$, $S_3=15\ksi$, 2024-T3 clad:
Case~1 $W=2.0/m=0.4$; Case~2 $W=4.0/m=1.5$ (inches throughout).

\subsection{Items that reconcile}

\begin{itemize}
\item \textbf{Tension solution}: identical. At the first step of Case~1 the
  benchmark reports $\beta^{S}_1=3.138$, $\beta^{S}_2=1.902$ and tension-only
  $\Delta K_1=8.344\ksqin$; this implementation gives $3.138/1.907$ and
  $8.34\ksqin$.
\item \textbf{Tension-only life}: benchmark $12{,}482$ vs.\ $12{,}602$ cycles
  (Case~1, $S_3=0$, our PZC off) --- within $1\%$. Material representation and
  integration are therefore equivalent for practical purposes.
\item \textbf{Plastic zone treatment}: the benchmark applies its plasticity
  logic only near instability (\texttt{IZPC} entries at the end of its
  output), whereas our optional Irwin correction inflates $K$ at every cycle
  ($\approx 9\%$ life reduction in these cases). Like-for-like comparisons
  should be run with our PZC off.
\item \textbf{Termination}: the benchmark's ``net section failure --
  ligament'' at $c_2\approx0.25$ vs.\ our link-up at $c_2=0.30$ differ by
  $<10$ cycles; growth through the last ligament fraction is effectively
  instantaneous and the criterion choice is immaterial here.
\end{itemize}

\subsection{The bearing-load discrepancy}
\label{sec:bearing}

The entire life discrepancy traces to the bearing (pin-load) $\beta$. The
benchmark's bearing betas (its BP columns) are far below the NASGRO TC23
Solution~C values at short cracks, and below the NASGRO TC05 FEM pin-load
tables as well (Table~\ref{tab:bp}).

\begin{table}[h]
\centering
\caption{Case 1 bearing $\beta$ (factor on $S_3\sqrt{\pi c_i}$) at the
initial crack state $c_1=0.01$, $c_2=0.05$.}
\label{tab:bp}
\begin{tabular}{lcc}
\toprule
Model & tip 1 & tip 2\\
\midrule
NASGRO TC23 Solution C            & 1.144 & 0.554\\
NASGRO TC05 FEM table (C7)        & 0.813 & 0.452\\
Effective width $W_{\rm eff}=7D$  & 0.450 & 0.272\\
Benchmark                          & 0.451 & 0.300\\
\bottomrule
\end{tabular}
\end{table}

Three observations from the benchmark's own output:
\begin{enumerate}
\item The ratio of its bearing to tension betas is nearly constant at the
  same value in both cases despite the width doubling:
  $\mathrm{BP}_1/\mathrm{BS}_1 \approx 0.143\ldots0.149 \approx 1/7$.
  This is the signature of an \emph{effective-width} bearing model,
  \begin{equation}
    K_{\rm bear} = \frac{D}{W_{\rm eff}}\,\beta^A\,\bar S_3\sqrt{\pi c},
    \qquad W_{\rm eff}=\min(7D,\,W),
    \label{eq:effw}
  \end{equation}
  i.e.\ the transferred load reacted locally over a few diameters
  (consistent with the $\sim35^\circ$ Inglis load-spread argument and the
  AFGROW guidance that bearing solutions should not see edges farther than
  $\sim6D$).
\item Its tip-2 ratio decays with crack growth ($0.141\to0.123$ over the
  Case~2 extract), i.e.\ its bearing beta lacks the long-crack
  edge-proximity magnification of the full-width $\beta^A$ --- suggesting the
  beta is evaluated in the $W_{\rm eff}$ sub-geometry rather than scaled
  from the full-width tension beta.
\item Eq.~\eqref{eq:effw} implemented directly reproduces the benchmark's
  BP$_1$ within $1$--$6\%$ at all extracted points in both cases, and lives
  within $+4\%$ (Case~1) / $-15\%$ (Case~2), Table~\ref{tab:lives}.
\end{enumerate}

NASGRO's Solution~C models the pin load as a point wedge force on the crack
faces (the $W/\pi c_0$ term in C2:3); this is asymptotically correct for
$c\gtrsim 2D$ but conservative at short cracks where the distributed
bore-pressure field decays rapidly --- the same conservatism documented for
analytical bearing betas vs.\ FEM in the AFGROW 2016 workshop (factor
$\approx3$ at $c=0.2D$). Three selectable bearing models are therefore
provided:

\begin{enumerate}
\item \textbf{NASGRO Solution C} (default; most conservative);
\item \textbf{Hybrid TC05$\to$TC23}: NASGRO TC05 FEM pin-load tables
  (Table~C7, strip$\leftrightarrow$row equivalence $D/H=D/W$,
  $x=c/(W-D)$) at short cracks, tapering linearly in $c/D$ to Solution~C
  over $c/D=0.5\ldots1.0$;
\item \textbf{Effective width}, Eq.~\eqref{eq:effw}, $k_{\rm brg}=7$
  (user-overridable) --- closest to the benchmark.
\end{enumerate}

\begin{table}[h]
\centering
\caption{Predicted lives (cycles), PZC off, vs.\ benchmark.}
\label{tab:lives}
\begin{tabular}{lcc}
\toprule
 & Case 1 ($W{=}2$, $m{=}0.4$) & Case 2 ($W{=}4$, $m{=}1.5$)\\
\midrule
No bearing ($S_3=0$)        & 12{,}602 & 33{,}206\\
NASGRO Solution C           &  5{,}524 & 20{,}493\\
Hybrid TC05$\to$TC23        &  6{,}588 & 20{,}263\\
Effective width $7D$        &  7{,}917 & 20{,}766\\
\textbf{Benchmark}          &  \textbf{7{,}594} & \textbf{24{,}508}\\
\bottomrule
\end{tabular}
\end{table}

\subsection{Questions for the benchmark author}

\begin{enumerate}
\item What is the basis of the bearing $\beta$ solution --- an effective-width
  construction per Eq.~\eqref{eq:effw}, a dedicated FEM table for the
  offset-hole strip, or other? Is the factor $1/7$ explicit?
\item Is the bearing beta evaluated in the $W_{\rm eff}$ sub-geometry (which
  would explain the declining tip-2 ratio at longer cracks), and how is the
  near free edge treated when $m<W_{\rm eff}/2$ (Case~1 shows the bearing
  beta turning back up near link-up, BP$_2=0.374$ at $c_2=0.25$)?
\item How is the pin load reacted in the bearing sub-problem (one end, as in
  NASGRO, or symmetrically)? The conventions differ by an open-hole tension
  case at $\tfrac12(D/W)S_3$ and require different tension inputs for the
  same physical joint.
\item Confirmation that ``Bearing Cut Off'' does not alter the bearing
  $\beta$ (verified empirically: identical life with the flag cleared).
\item At short cracks the benchmark's betas sit below the NASGRO TC05 FEM
  family ($0.30$ vs.\ $0.41$--$0.45$ at $c_2=0.05$); what validation
  supports the lower values?
\end{enumerate}

\section{Skin bridging (TC23B)}

\subsection{Displacement compatibility}

Bridging follows Swift's displacement-compatibility method. By symmetry the
fasteners act in pairs at $y_j=\pm jp$, $j=1\ldots n$ (the fastener in the
cracked hole lies on the crack plane and transfers no bridging load). With
pair forces $F_j$ and fastener flexibility $f$, compatibility of the
relative stringer--skin displacement at each fastener requires
\begin{equation}
  \bigl[\,f\,\mathbf{I} + \mathbf{G}^{\rm str} + \mathbf{G}^{\rm sk}\,\bigr]\,
  \mathbf{F} \;=\; \mathbf{V}_0\,\sigma \;+\; \Delta\varepsilon\,\mathbf{y},
  \label{eq:compat}
\end{equation}
solved at each integration step for the current effective half-crack
$a=c_0=(c_1+D+c_2)/2$ (centre-crack idealisation of hole plus cracks).

\paragraph{Opening term.} $V_0(y_j)$ is the crack-induced opening of the
stringer at the fastener under unit remote stress, from the Westergaard
field (plane stress):
\begin{equation}
  V_0(x,y)=\frac1{E_1}\Bigl[\,2\,\mathrm{Im}\bigl(\sqrt{z^2-a^2}-z\bigr)
  -(1+\nu_1)\,y\,\Bigl(\mathrm{Re}\tfrac{z}{\sqrt{z^2-a^2}}-1\Bigr)\Bigr],
  \quad z=x+\mathrm{i}y .
\end{equation}

\paragraph{Skin influence.} $G^{\rm sk}_{ij}$ is the displacement of the
intact infinite skin at fastener $i$ due to the reacted pair at $j$,
evaluated from the exact Muskhelishvili point-force (Kelvin) potentials
\begin{equation}
  \varphi(z)=A\ln(z-z_0),\quad
  \psi(z)=-\kappa\bar A\ln(z-z_0)-\bar z_0 A/(z-z_0),\quad
  A=-\tfrac{f_x+\mathrm{i}f_y}{2\pi(1+\kappa)},
\end{equation}
with $\kappa=(3-\nu)/(1+\nu)$. The self-term log singularity is regularised
by averaging over a circle of radius $d/2$ (deformation inside the cutoff is
represented by the empirical fastener flexibility).

\paragraph{Cracked-sheet influence.} $G^{\rm str}_{ij}$ adds to the Kelvin
field a crack-release correction obtained \emph{exactly} from the
Betti/Rice weight-function identity, which collapses to a compliance
integral of the pair SIF Green's functions:
\begin{equation}
  v^{\rm corr}_{ij}
  = -\frac{t_1}{E_1}\int_0^{a}
     \bigl[K^+_i(a')K^+_j(a')+K^-_i(a')K^-_j(a')\bigr]\,\mathrm{d}a' ,
  \label{eq:betti}
\end{equation}
where $K^{\pm}_j(a')$ is the centre-crack SIF at the $\pm a'$ tip due to a
unit restraining pair at $y_j$, by Bueckner superposition of the Kelvin
crack-line stress through the centre-crack weight functions:
\begin{equation}
  K^{\pm}_j(a')=\frac{1}{\sqrt{\pi a'}}\int_{-a'}^{a'}
  \sigma^{\rm pair}_{yy}(x,0)\,\frac{a'\pm x}{\sqrt{a'^2-x^2}}\,\mathrm{d}x
  \qquad\text{(Gauss--Chebyshev quadrature).}
\end{equation}
No fitted influence functions are used anywhere in the chain.

\paragraph{Biaxial skin stress.} An off-axis biaxial skin state alters the
stringer/skin load-sharing through Poisson contraction. Under proportional
loading (skin stresses $\sigma_L,\sigma_H$ specified at peak stringer stress
$\hat\sigma$ and scaling with it), the far-field strain mismatch per unit
stringer stress is
\begin{equation}
  \Delta\varepsilon/\sigma=\frac1{E_1}
  -\frac{\sigma_L-\nu_2\,\sigma_H}{E_2\,\hat\sigma},
\end{equation}
entering Eq.~\eqref{eq:compat} as the relative slip $\Delta\varepsilon\,y_i$
measured from the crack plane. Strain-matched skin ($\Delta\varepsilon=0$) is
the default. As noted in \S\ref{sec:bypass}, the mismatch acts mainly as a
no-crack load-share (subtracted from $P_{\rm tot}$) and so has only a
second-order effect on crack-growth life under the bypass route.

\subsection{SIF reduction --- load bypass}
\label{sec:bypass}

The solved pair forces reduce the tip SIF through a \emph{load-bypass}
statement: the load shed into the skin bypasses the crack section, so the
membrane stress reaching the crack is the bypass stress
\begin{equation}
  P_{\rm app}=S_{\rm gross}\,W\,t,\quad
  P_{\rm tot}=\sum_j F_j(a)-\sum_j F_j(0),\quad
  R_{\rm bp}=\frac{S_{\rm byp}}{S_{\rm gross}}
            =1-\frac{P_{\rm tot}}{S_{\rm gross}\,W\,t},
\end{equation}
with $S_{\rm gross}=\sigma+(D/W)S_3$, applied to the membrane (tension)
contribution of $K$ only:
\begin{equation}
  K^{\rm TC23B}=R_{\rm bp}\,K_{\rm tension}+K_{\rm bending}+K_{\rm bearing},
\end{equation}
since in-plane bending ($S_2$) and pin-bearing ($S_3$) are separate load
paths that do not bypass. $P_{\rm tot}$ is the \emph{crack-driven} one-sided
bypassed load ($F_j$ = upper fastener of pair $j$; the lower returns it below
the crack). The no-crack baseline $\sum_j F_j(0)$ is subtracted because the
far-field strain mismatch (\S\,biaxial) produces distributed load-sharing that
exists with no crack and does not bypass it; it also removes a
finite-fastener-row artefact in which the raw $\sum_j F_j$ grows with the
modelled row length. The baseline is identically zero in strain-matched mode.
$R_{\rm bp}>1$ (load attraction) is admitted when the skin is strained more
than the stringer ($R_{\rm bp}$ clamped to $[0.05,2]$). Because
$F_j\propto\sigma$ under proportional loading, $R_{\rm bp}$ is
stress-independent and the constant-amplitude engine is unchanged.

This replaces the weight-function crack-face-traction reduction
$R=\bigl(\sqrt{\pi a}+\sum_j F_j K^+_j(a)\bigr)/\sqrt{\pi a}$ (retained as a
diagnostic). The two routes account for the \emph{same} fastener loads --- the
bypass route is the global load-path statement, the weight-function route the
local crack-tip-shielding statement --- so applying both would double-count.
The compatibility solve itself is unchanged: it uses the rigorous
crack-release compliance, Eq.~\eqref{eq:betti}, to obtain the correct
equilibrium $F_j$; only the $K$-from-$F_j$ step changed.

A consequence is that the far-field strain mismatch (\S\,biaxial) becomes a
\emph{minor} influence on crack-growth life under the bypass route: its primary
action is a uniform stringer/skin load-share that exists with no crack (the
subtracted baseline) and is already represented in the input stringer stress;
only the small crack--mismatch coupling remains. To make a hoop/biaxial
load-share retard the crack as a first-order effect it must instead be applied
as a reduction of the input far-field stringer stress (a stiffness-ratio
load-sharing calculation outside this tool's scope), not through crack bypass.

In the post-link-up SENT phase the edge crack of depth $a_e$ is modelled as
a half centre-crack with the model centre at the free edge, the fastener
line at depth $m$, and free-edge image sources at $-m$; the skin (which has
no free edge) sees the real forces only.

\subsection{Fastener flexibility --- modified Tate \& Rosenfeld}

\begin{equation}
  f=\underbrace{\frac{0.375}{E_f t_1}+\frac{0.375}{E_f t_2}}_{\text{fastener bearing}}
  +\underbrace{\frac{0.9}{E_1 t_1}+\frac{0.9}{E_2 t_2}}_{\text{sheet hole compliance}}
  +\frac{32(1+\nu_f)(t_1+t_2)}{9\,E_f\,\pi d^2}
  +\frac{8\,(t_1^3+5t_1^2t_2+5t_1t_2^2+t_2^3)}{5\,E_f\,\pi d^4},
\end{equation}
in $\mathrm{in/kip}$ with $E$ in ksi (sheet 1 = stringer, sheet 2 = skin).

\subsection{Outputs and checks}

Per-fastener bridging loads $F_j\sigma$ are reported at the final crack size
together with the peak over the run, and checked against an optional shear
allowable. Fastener failure is \emph{not} progressively released; results
beyond first exceedance are flagged as unconservative. Net-section yield
ignores the load shed to the skin (conservative). With a biaxial mismatch
the outermost fastener loads represent joint end transfer and depend on the
modelled row length $n$; the near-crack loads and the SIF restraint do not.

\section{Verification}

All checks run headless (\texttt{node tests/run-tests.mjs}; 36 assertions):

\begin{itemize}
\item \textbf{Wedge-force limit}: $K^{\pm}$ of a restraining pair
  $\to -P/(t\sqrt{\pi a})$ as the pair height $\to0$ (Richardson-extrapolated,
  $2\%$).
\item \textbf{Equilibrium}: $\int\sigma^{\rm pair}_{yy}(x,0)\,\mathrm{d}x=-1/t$
  for the half-plane containing one force of the pair.
\item \textbf{Field consistency}: $\partial v/\partial y=
  (\sigma_{yy}-\nu\sigma_{xx})/E$ for the Kelvin displacement vs.\ stress
  fields.
\item \textbf{Independent cross-check}: the closed-form Westergaard $V_0$
  equals the Betti route
  $-(t_1/E_1)\int\sqrt{\pi a'}\,(K^++K^-)\,\mathrm{d}a'$ built from the
  Kelvin stresses and weight functions ($0.5\%$) --- validating the entire
  Green's-function pipeline.
\item \textbf{Reciprocity}: the assembled influence matrix is symmetric
  (Betti) to $10^{-6}$.
\item \textbf{Regression}: bridging off (or $f\to\infty$) reproduces plain
  TC23 exactly; the bypass factor equals $1-P_{\rm tot}/P_{\rm app}$, is
  stress-independent, and reduces the membrane term only; bypass deepens with
  crack length; bridged life exceeds unbridged; the no-crack bypass factor is
  $\approx1$ in every skin-stress mode (baseline correctly removed); biaxial
  life stays within a few percent of strain-matched life; the modified
  Tate \& Rosenfeld value is pinned by hand calculation.
\item \textbf{Benchmark}: tension-only life within $1\%$; bearing betas and
  lives per \S\ref{sec:bench}.
\end{itemize}

\section{Assumptions and limitations}

\begin{enumerate}
\item Bridging influence functions use infinite-sheet elasticity for the
  stringer; finite width enters only through the TC23 $\beta$ factors
  (standard compounding). Plane stress throughout.
\item The hole is omitted from the bridging displacement field (centre-crack
  idealisation $2c_0=c_1+D+c_2$); the fastener line is taken through the
  flaw centre (exact for $c_1=c_2$).
\item The restraint ratio is identical at both tips in the two-crack phase.
\item Constant amplitude, proportional loading; no load-interaction or
  progressive fastener failure modelling.
\item The SENT-phase half-model approximates the free edge by image sources;
  the bridging correction there carries the usual edge-crack/centre-crack
  approximation.
\end{enumerate}

\begin{thebibliography}{9}
\bibitem{nasgro} NASGRO Reference Manual v8.2, Southwest Research
  Institute --- Appendix~C, crack cases TC23 (pp.~C-36--43) and TC05.
\bibitem{swift84} T.~Swift, ``Fracture Analysis of Stiffened Structure,''
  in \emph{Damage Tolerance of Metallic Structures}, ASTM STP 842, 1984.
\bibitem{swift90} T.~Swift, ``Repairs to Damage Tolerant Aircraft,''
  FAA-AIR-90-01, 1990.
\bibitem{tate} M.~B. Tate and S.~J. Rosenfeld, ``Preliminary Investigation
  of the Loads Carried by Individual Bolts in Bolted Joints,'' NACA TN-1051,
  1946 (modified form with 0.375/0.9 hole-compliance numerators as used in
  industry practice).
\bibitem{musk} N.~I. Muskhelishvili, \emph{Some Basic Problems of the
  Mathematical Theory of Elasticity}, Noordhoff (point-force potentials).
\bibitem{rice} J.~R. Rice, ``Some Remarks on Elastic Crack-Tip Stress
  Fields,'' \emph{Int.\ J. Solids Structures}, 8, 1972 (weight-function
  relations).
\bibitem{harter} J.~A. Harter and A.~V. Litvinov, ``Modeling Bearing Load in
  Wide Panels Using AFGROW,'' AFGROW European Workshop, 2016 (analytical
  vs.\ FEM bearing betas; effective-width guidance).
\bibitem{tada} H.~Tada, P.~C. Paris and G.~R. Irwin, \emph{The Stress
  Analysis of Cracks Handbook}, 3rd ed., ASME Press, 2000.
\end{thebibliography}

\end{document}
